\begin{abstract}
    Recent work has explored the potential of and graph neural networks (GNNs) in particular, to tackle combinatorial problems on graphs. While many such problems are NP-hard and therefore unlikely to admit efficient algorithms in the worst case, it remains plausible that GNNs could nevertheless succeed on instances of practical relevance. In this work, we examine this question via the \emph{clique--coloring problem} (CPP), a polynomial-time solvable promise problem distinguishing graphs with a $k$-clique from those that are $(k\!-\!1)$-colorable.

    We prove that standard message-passing GNNs are fundamentally unable to distinguish these two classes of graphs. For every $k \ge 3$, we construct an infinite family of graph pairs, one containing a $k$-clique and the other $(k\!-\!1)$-colorable, that are indistinguishable by any such GNN. This result is unconditional and does not rely on complexity-theoretic assumptions. We further complement our theoretical findings with experiments, which illustrate that while CPP is easy for GNNs on planted random graph distributions, performance deteriorates when the planted $k$-clique and $(k\!-\!1)$-colorable distributions are adjusted to have matching average degree.

\end{abstract}